\section[Local multipartite core responses]{Local discovery and incremental multipartite core responses}
\label{sec:op3-local-multipartite}

\paragraph{Status and model.}
This \statusproved{} draft, awaiting independent review, extends the scalar
clique construction to an unknown complete multipartite core. The complete
local state machine and incremental event heap are implemented. Neither a
partition nor a core size nor a future support is supplied. The input graph
is promised to consist of a complete multipartite core $C$, with at least
two nonempty parts $C_g$, and $t_i\geq0$ private degree-one leaves at each
core vertex $i$. Put $N=|C|$ and $n_g=|C_g|$. For this theorem assume
\begin{equation}
 N-\max_g n_g\geq2,\qquad v\in C.
 \label{eq:op3-multipartite-promise}
\end{equation}
Thus original degree distinguishes every core vertex from every pendant.
The promise is an assumption; it is not globally verified on unread rows.
Use finite simple connected unweighted graphs, original full degrees,
local degree/adjacency access, $0<\alpha<1$, $\lambda>0$, and exact-real
word arithmetic. The physical response and original residual conventions
are unchanged. This is a structural result; arbitrary-graph OP3 remains
\statusopen{}.

\begin{lemma}[One scalar coordinate for all discovered parts; \statusproved]
\label{lem:op3-multipartite-scalar}
Suppose the identifiers and original degrees of all core vertices are
known, and let $\mathcal J$ be a collection of identified complete parts
containing the seed part. Set every other core coordinate to zero and
solve the obstacle problem on the identified parts and their private
leaves. Write $S=\sum_{i\in C}u_i$ and
$S_g=\sum_{i\in C_g}u_i$. For $i\in C_g$, put
$b_i=\mathbb 1_{\{i=v\}}-\lambda d_i$ and
\[
 \psi_i(T)=
 \begin{cases}
 0,&b_i+T\leq0,\\
 (b_i+T)/d_i,&0<b_i+T\leq d_i\lambda/\gamma,\\
 (b_i+T-\gamma\lambda t_i)/(d_i-\gamma^2t_i),
     &b_i+T>d_i\lambda/\gamma.
 \end{cases}
\]
Let $F_g(T)=\sum_{i\in C_g}\psi_i(T)$ and define $G_g(S)$ by the
unique solution of $T/\gamma+F_g(T)=S$, with $G_g(S)=F_g(T)$.
Then the restricted obstacle solution is determined by the unique
nonnegative scalar root
\begin{equation}
 S=\sum_{g\in\mathcal J}G_g(S).
 \label{eq:op3-multipartite-root}
\end{equation}
On a piece $F_g(T)=AT+B$,
\begin{equation}
 G_g(S)=\frac{\gamma AS+B}{1+\gamma A},\qquad
 S_* = T_*/\gamma+F_g(T_*)
 \label{eq:op3-part-transform}
\end{equation}
is the corresponding transformed piece and breakpoint. Part $g$ has at
most $2n_g$ such events, with
\begin{equation}
 0\leq G'_g\leq
 \frac{\gamma n_g}{N-n_g+\gamma n_g}<\frac{n_g}{N}.
 \label{eq:op3-part-slope}
\end{equation}
\end{lemma}
\begin{proof}
Each private leaf has value $(\gamma u_i-\lambda)_+$. Substitution into
the original core equation gives
\[
 d_i u_i-\gamma t_i(\gamma u_i-\lambda)_+
       =b_i+\gamma(S-S_g)
\]
at a positive coordinate, with its nonpositive original gate at zero.
The two positive denominators obey
$d_i-\gamma^2t_i=N-n_g+(1-\gamma^2)t_i>0$, so solving the pieces gives
$u_i=\psi_i(\gamma(S-S_g))$. The part equation is exactly the displayed
inverse relation. Its left side is continuous, strictly increasing and
onto the real line. Negative $T$ intervals must be retained when building
the inverse: for trial $S$ below the final root the seed part can have
$G_g(S)>S$. At the final root, all part masses are nonnegative and sum to
$S$, so every actual field $\gamma(S-S_g)$ is nonnegative.

The affine inverse gives \cref{eq:op3-part-transform}; continuity makes
the breakpoint value independent of which adjacent piece is used.
The transformation is strictly increasing, preserving event order.
Every slope of $F_g$ is at most $n_g/(N-n_g)$, giving
\cref{eq:op3-part-slope}. Summing over identified parts gives a slope
strictly less than one, even if only some parts are present.
The sum of part responses is nonnegative at zero; subtracting $S$ is
strictly decreasing and eventually negative. Its unique nonnegative root
satisfies all original restricted obstacle conditions. Positive definiteness
identifies it as the unique restricted optimum.
\end{proof}

\begin{lemma}[Two positive rows pay for complete core identifiers; \statusproved]
\label{lem:op3-multipartite-startup}
Either the exact obstacle solution is zero or a seed star, or two
support-safe original row scans identify all core labels and the seed's
and one adjacent core vertex's entire parts. All startup work is
$O(1+(d_v+d_j)\log(2+d_v+d_j))$, including degree queries, dictionaries
and copies; the row $j$ is read only in the latter case.
\end{lemma}
\begin{proof}
If $\lambda d_v\geq1$, all physical loads are nonpositive and zero is
optimal. Otherwise $v$ is positive. Read its row and query every neighbor's
original degree. The degree-one neighbors are precisely its $t_v$ private
leaves; the others are all core vertices outside the seed's part.
The exact restricted seed-star value is
$u_v=\psi_v(0)$, using the now known $d_v,t_v$. A degree-one leaf's
equation is determined by that incidence and its degree, so this virtual
star requires no additional leaf-row scan before final output.

For each exposed core vertex $j$, test the original excess
$\gamma u_v-\lambda d_j$. If all are nonpositive, every exposed core
gate is quiet. Other vertices in the seed's part see only zero core
neighbors, and every remaining private leaf has zero parent. These are
the complete original obstacle conditions, proving the seed-star stop.

Otherwise choose a strictly positive gate and read that vertex $j$'s row.
The usual inverse-positive obstacle admission argument makes $j$ positive
in the full optimum, so this scan is support-safe. Since $v,j$ lie in
different parts, every core vertex is adjacent to at least one of them.
The union of the two rows' degree-greater-than-one neighbors therefore
contains all core labels, including $v,j$. Their number is at most
$d_v+d_j$, and all their degrees have been queried.

The seed part consists exactly of core labels not in the seed's neighbor
set, including $v$ itself; the analogous complement of $j$'s row is its
entire part. They are disjoint. Building these sets and core-label records
costs linear many dictionary operations and copies in the two row lengths.
Comparison dictionaries give the stated bound. Both scanned rows are cached
for final output and later classification.
\end{proof}

\paragraph{Incremental part and event state.}
After the two-row startup, build the seed-part curve and solve its scalar
root, then insert the second identified part and solve again. Each new
part's size and its known original degrees give $t_i=d_i-(N-n_g)$.
Build its at most $2n_g$ original breakpoints, sort them, and transform
them by \cref{eq:op3-part-transform}. At the current exact scalar root,
consume all past transformed events to initialize this part's current
affine piece. Insert only its future events in one global min-heap.
Store one current affine coefficient pair per part, and one aggregate
coefficient pair for their sum. Earlier parts' curves are not rebuilt.

On a current combined interval write $\sum_gG_g(S)=AS+B$. Compute the
candidate $B/(1-A)$. If it lies at or before the next event, it is the exact
restricted root. Otherwise cross that event, update its part's coefficient
pair and the aggregate pair, and continue. A provisional affine root is
never used as an original boundary certificate. Equality at an event is
valid in both adjacent continuous pieces.

Keep unidentified core labels in a dictionary $R$ and a second min-heap
with keys $\lambda d_i/\gamma$. Stale entries from identified parts are
removed with charged work. At each exact restricted root, either
$S\leq\min_{i\in R}\lambda d_i/\gamma$ and stop, or the heap names a
strictly positive original gate $j$. Read its row, classify the current
$R$ by adjacency to $j$, remove the nonneighbors as its newly identified
part, build that part once, and continue. Original degrees and penalties
are never replaced by induced degrees.

\begin{lemma}[Paid discovery and monotone event consumption; \statusproved]
\label{lem:op3-multipartite-paid-events}
Every representative scanned by this continuation is positive in the
full optimum. Total complement-classification work uses
$O(N+\sum_jd_j)$ dictionary operations and copied labels, where the sum
runs over the newly scanned representatives. Across all identified parts,
the algorithm creates at most $2N$ scalar events and consumes each at most
once, with $O(N\log(2+N))$ total curve and heap work.
\end{lemma}
\begin{proof}
Every unidentified vertex is adjacent to every identified core vertex.
Thus its true original gate at the restricted optimum is
$\gamma S-\lambda d_j$, regardless of unequal values within identified
parts. The heap minimum is an exact quietness test. A positive gate gives
a support-safe representative before its row is read.

Among currently unidentified labels, the nonneighbors of $j$, including
$j$ itself, are exactly its complete part. During a scan of $R$, each
nonneighbor is removed permanently and charged once to its identifier.
Every other inspected label is an actual neighbor of $j$ and can be charged
to that newly scanned original row. The same accounting pays for a temporary
snapshot of $R$ used while deleting entries; it is not a free complement
oracle. All unknown-heap insertions and stale removals number at most $N$.

Adding an identified part enlarges the virtual obstacle region, so its
restricted solution increases entrywise. Equivalently, its nonnegative
part response increases the right side of \cref{eq:op3-multipartite-root}.
For a positive representative the increase is strict at the previous root.
The new root is therefore no smaller, and all previously crossed events
remain in the past. Each newly identified vertex contributes at most two
events, built once. Its past events are consumed during initialization;
its future events are inserted and subsequently crossed at most once.
The algorithm keeps its exact scalar root until the next part insertion,
so every root search costs one final test plus its newly crossed events.
Sorting separately within parts costs
$O(\sum_g n_g\log(2+n_g))\leq O(N\log(2+N))$.
Heap operations and constant-size coefficient updates obey the same bound.
No old coordinate-value vector is recovered between part insertions.
\end{proof}

\begin{theorem}[Local exact multipartite-core solver; \statusproved]
\label{thm:op3-local-multipartite}
On the promised family \cref{eq:op3-multipartite-promise}, the preceding
algorithm returns the exact $\lambda$-obstacle response. If $U$ is its
returned positive support, total local exact-real word work and space are
\begin{equation}
 O\!\left(1+\operatorname{cvol}(U)
                    \log(2+\operatorname{cvol}(U))\right),
 \qquad O(1+\operatorname{cvol}(U)),
 \label{eq:op3-multipartite-local-work}
\end{equation}
respectively, including all graph accesses, dictionary operations, copies,
events and output. For $0<\epsappr,\rho<1$, choosing
$\lambda=\epsappr/2$ gives $\widetilde O(1/\epsappr)$ ACL work;
separately, choosing $\lambda=\rho$ gives exact
$\widetilde O(1/\rho)$ RPPR work on this family.
\end{theorem}
\begin{proof}
The startup lemma handles both early exits. In the remaining case,
\cref{lem:op3-multipartite-scalar,lem:op3-multipartite-paid-events} give
exact restricted solutions and support-safe new representatives until all
unidentified gates are nonpositive. The identified-part obstacle equations,
unknown-core inequalities and their zero-parent leaf inequalities are the
full original obstacle conditions, so the terminal restricted solution is
the unique full optimum.

At termination evaluate each identified part's mass and then each of its
core values once. Read a core row only when its exact value is positive,
reusing cached representative rows. Its private leaves have the known
common value $(\gamma u_i-\lambda)_+$; enumerate them along that positive
row and emit positive coordinates. If desired, as in the implementation,
read each positive degree-one row once as an incidence check. Never read a
zero core or zero leaf row. Unidentified core values are implicitly zero.
These final operations cost $O(\operatorname{vol}(U))$ dictionary operations
and original incidence reads, including reads of cached rows.

Both startup representatives and every later representative lie in $U$.
Consequently $N\leq d_v+d_j\leq\operatorname{vol}(U)$ after startup,
and the sum of all representative degrees is at most
$\operatorname{vol}(U)$. This pays for all core metadata and even the
identified core values that turn out to be zero. The preceding lemmas
and comparison dictionaries prove the work bound. Current event heaps,
part records, degree/boundary dictionaries, cached positive rows and output
contain $O(1+\operatorname{cvol}(U))$ words; temporary part curves and
complement snapshots have the same peak bound and their construction is
charged. No old complete response histories are retained.

The original obstacle mass identity gives
$\lambda\operatorname{vol}(U)<1$ for nonzero output, yielding the two
separate accuracy consequences. There is no polynomial arithmetic-operation
dependence on $1/\alpha$, no global preprocessing, and no imported SDD or
dynamic-tree solver in this construction. Bit complexity, unchecked
floating-point stability and an arbitrary-core theorem are not asserted.
\end{proof}

\paragraph{Measured and open extensions.}
The actual state machine in \nolinkurl{local_multipartite_pendants.py}
is compared against independent full original obstacle solutions, including
restricted optima between part insertions and original unknown-coordinate
gates. Separate ledgers check that old events are retained, every coordinate
curve is built once, and complement scans are paid. Reversed incidence order,
exact first-core/leaf/unknown-part/zero gate ties, and enormous implicit
inactive attachments are included in
\nolinkurl{LOCAL_MULTIPARTITE_PENDANT_AUDIT.json}.
All dense matrices are validators only. The Python driver records abstract
dictionary operations; the deterministic word theorem prices them using
comparison dictionaries, rather than assuming worst-case constant-time
Python hashing. The next subsection removes
\cref{eq:op3-multipartite-promise}; general OP3 and recursive module
discovery remain separate \statusopen{} targets.

\subsection{Canonical core and arbitrary physical seed}
\label{sec:op3-canonical-multipartite}

\begin{lemma}[The non-leaf core is canonical; \statusproved]
\label{lem:op3-canonical-core}
On the same promised graph family, without
\cref{eq:op3-multipartite-promise}, let
$C^\circ=\{i:d_i>1\}$. If $|C^\circ|\geq2$, its induced graph is a
connected complete multipartite graph, and every other vertex is its
private leaf. If $|C^\circ|=1$, the graph is a star. If
$C^\circ=\varnothing$, it is the two-vertex graph. Thus original degree
provides the required classification whenever a positive core row is read;
no canonical partition or global classification is supplied.
\end{lemma}
\begin{proof}
Every attached private leaf has degree one. An original core vertex in
part $g$ can have degree one only if $N-n_g=1$ and it has no attached
leaves. When $N>2$, this forces the original core to be a star: one part
contains its center, and the other contains its peripheral vertices.
Removing peripheral core vertices of original degree one leaves a star
core, possibly just its center. Each removed vertex is now a private leaf
of that center. When $N=2$, keeping both endpoints, one endpoint or neither
gives respectively an edge core, a singleton core or the two-vertex graph.
In all other cases every original core vertex has degree greater than one,
so the original multipartite representation is already canonical.
The conclusions concern the promised graph; they do not claim to verify
the promise using unread adjacency rows.
\end{proof}

\begin{lemma}[Eliminating a physical leaf seed; \statusproved]
\label{lem:op3-multipartite-leaf-seed}
Suppose $d_v=1$ and $\lambda<1$, and let $a$ be the neighbor of the
physical seed $v$. The seed-only solution is exact if
\begin{equation}
 \gamma(1-\lambda)-\lambda d_a\leq0.
 \label{eq:op3-leaf-anchor-gate}
\end{equation}
Otherwise $a$ is positive in the full obstacle solution and its row may
be read. Eliminating the necessarily positive seed gives
\begin{equation}
 u_v=1-\lambda+\gamma u_a,\qquad
 \widehat d_a=d_a-\gamma^2,\qquad
 \widehat b_a=\gamma(1-\lambda)-\lambda d_a.
 \label{eq:op3-forced-seed-kernel}
\end{equation}
All other loads and diagonal coefficients are unchanged, and the seed is
excluded from the anchor's ordinary thresholded leaves. Every penalty
still uses the original full degree.
\end{lemma}
\begin{proof}
Since $\lambda<1$, the original seed equation has positive right side
$1-\lambda+\gamma u_a$, so the displayed recovery holds for every
nonnegative neighbor value. With $u_a=0$, the only possible violated
boundary inequality is exactly \cref{eq:op3-leaf-anchor-gate}. A positive
gate certifies a positive anchor by obstacle comparison before its row is
scanned. Substituting the seed equation into the anchor equation gives
\cref{eq:op3-forced-seed-kernel}. This is an exact elimination from the
original positive definite problem. It changes the equation coefficients,
not the original degree in $-\lambda d_a$ or in boundary certificates.
\end{proof}

\paragraph{Modified coordinate kernels.}
If the physical seed is a core vertex, take it as anchor and use the
previous kernels. Otherwise first perform the leaf-seed test and, on a
positive gate, use the forced-seed elimination above. For a canonical core
vertex $i$, let $c_i$ be its canonical core degree, let $\sigma_i$ indicate
that it is the forced seed's anchor, and set
\[
 a_i=d_i-\gamma^2\sigma_i,\qquad
 q_i=d_i-c_i-\sigma_i,\qquad
 b_i=
 \begin{cases}
 \gamma(1-\lambda)-\lambda d_i,&\sigma_i=1,\\
 \mathbb 1_{\{i=v\}}-\lambda d_i,&\sigma_i=0.
 \end{cases}
\]
The response $\psi_i(T)$ is obtained by replacing $d_i,t_i,b_i$ in
\cref{lem:op3-multipartite-scalar} by $a_i,q_i,b_i$, respectively.
These replacements affect only that three-piece coordinate formula;
all original gate and accuracy formulas retain $d_i$. In particular,
\begin{equation}
 a_i-\gamma^2q_i
 =c_i+(1-\gamma^2)(d_i-c_i)\geq c_i.
 \label{eq:op3-canonical-kernel-slope}
\end{equation}
Whenever the canonical core has at least two vertices, $c_i=N-n_g\geq1$.
The denominator is positive, the bound on $F'_g$ and hence
\cref{eq:op3-part-slope} continue to hold, and each coordinate still
has at most two events. Negative transformed intervals remain necessary.

\begin{theorem}[Every physical seed, without a minimum core degree; \statusproved]
\label{thm:op3-canonical-multipartite}
For every finite simple connected unweighted graph consisting of a
complete multipartite core with private leaves, and every physical seed,
the canonical-core algorithm returns the exact $\lambda$-obstacle response
using only original degree/adjacency queries. No core size, partition or
seed-type annotation is supplied. Its exact-real word work, space, and
separate ACL and RPPR consequences are those of
\cref{thm:op3-local-multipartite}.
\end{theorem}
\begin{proof}
The initial $\lambda d_v\geq1$ test again returns zero exactly. Otherwise
a degree-one physical seed is positive; read its one-entry row and test
\cref{eq:op3-leaf-anchor-gate}, querying only its neighbor's degree.
A quiet gate returns the seed-only answer. A positive gate certifies the
anchor and permits its row scan and forced-seed elimination. A physical
core seed is itself the positive anchor without this preliminary step.

Use original degree to split the anchor's row into canonical core neighbors
and ordinary leaves, retaining the physical seed as a distinguished leaf
when necessary. Solve its exact virtual star using the modified kernel and
test each exposed canonical core gate $\gamma u_a-\lambda d_j$.
If all are quiet, the star and forced-seed recovery satisfy all original
obstacle conditions. This also handles a singleton canonical core and the
two-vertex graph. Their star kernel has a positive denominator: even when
$c_a=0$, $d_a>0$ and $\gamma<1$ give positivity in
\cref{eq:op3-canonical-kernel-slope}.

If a core gate is positive, the canonical core has at least two vertices.
Two positive rows enumerate its identifiers exactly as before. Use the
modified kernels in the same incremental part/event construction, with
$c_i=N-n_g$ computed after identifying a part.
\Cref{eq:op3-canonical-kernel-slope} supplies every bound used in
\cref{lem:op3-multipartite-scalar,lem:op3-multipartite-paid-events}.
Every unknown core vertex has load $-\lambda d_i$ and is adjacent to all
identified core vertices; it is not adjacent to a private physical seed.
Consequently its original gate is still $\gamma S-\lambda d_i$.
The exact quietness and monotone part-insertion arguments are unchanged.

At final recovery, skip the distinguished physical seed in its anchor's
ordinary-leaf loop and emit it once using
\cref{eq:op3-forced-seed-kernel}. Cache its already read row rather than
reading it again. Every scanned row, including this preliminary seed row,
belongs to the returned positive support. Both startup core rows and each
subsequent representative still pay for all metadata, complement scans,
curve construction and final output. The extra seed record and its gate
test cost only constant many words and dictionary operations. This proves
the same work and space bound. Finally the \emph{original} obstacle mass
identity gives $\lambda\operatorname{vol}(U)<1$; the eliminated load is
never substituted into that accuracy argument.
\end{proof}

\paragraph{Measured scope.}
The implemented extension in
\nolinkurl{canonical_multipartite_pendants.py} shares the actual incremental
event engine with the narrow solver. Its full exact audit has 4,840 explicit
original obstacle comparisons, including 2,664 physical leaf seeds and
1,650 cases with an original degree-one core vertex, 4,348 independent
restricted-optimum checks, and 56 implicit original-equation certificates.
Twenty-four included exact-tie cases cover the source-leaf gate, an ordinary
leaf born after elimination, an unknown-part gate and the zero solution.
The audit counts original graph accesses separately from dense validation,
checks the distinguished seed's original recovery, and verifies that the
set of scanned rows equals the returned positive support. The result is
a promised-family theorem and a proof draft awaiting independent review;
it is not an arbitrary-graph OP3 theorem.
